% Part: proof-theory
% Chapter: normalization
% Section: permutations

\documentclass[../../../include/open-logic-section]{subfiles}

\begin{document}

\olfileid{pt}{nor}{per}

\olsection{क्रमपरिवर्तन रूपांतरणे}

~\Log{N1i} मधील कट
हे !!a{elimination} नियमाच्या
प्रमुख आधारविधानात संपणारे
खंड असतात. सामान्यरूपीकरण
सिद्धतेतील एक प्रसंग म्हणजे
महत्तम कटांची लांबी आणि
त्याबरोबर !!{proof} ची
कट-लांबी कमी करणे.
$>1$ लांबीचा कट
अनेक प्रकारे संपू शकतो:
कटमधील~$!A$ ची शेवटची
!!{formula} आस्थिती
(\Elim{\lor} किंवा
\Elim{\lexists}) चा
निष्कर्ष आहे का, आणि
ती कोणत्या !!{elimination}
अनुमानाचे प्रमुख आधारविधान
आहे, यांवर ते अवलंबून असते.

उदाहरणार्थ, संबंधित
अनुमाने~\Elim{\lor} आणि
\Elim{\land} असतील,
$!A \ident !B \land !C$
असेल, आणि~$\delta_1$
ही उप-!!{proof} पुढील
रूपात संपत असेल:
\[
\AxiomC{}
\RightLabel{$\delta_2$}
\DeduceC{$!D \lor !E$}
\AxiomC{$\Discharge{!D}{x}$}
\RightLabel{$\delta_3$}
\DeduceC{$!B \land !C$}
\AxiomC{$\Discharge{!E}{y}$}
\RightLabel{$\delta_4$}
\DeduceC{$!B \land !C$}
\DischargeRule{\Elim{\lor}}{x\,y}
\TrinaryInfC{$!B \land !C$}
\RightLabel{\Elim{\land}[1]}
\UnaryInfC{$!B$}
\DisplayProof
\]
तर या कटमधील शेवटची
!!{formula} आस्थिती~$!A_n$
ही~\Elim{\lor} चा
निष्कर्ष, म्हणजे खालचे
$!B \land !C$, आहे.
त्याआधीची !!{formula}
आस्थिती~$!A_{n-1}$
ही~\Elim{\lor}
च्या गौण आधारविधानांपैकी
एक आहे.

या कटाची लांबी
कमी करण्यासाठी
\Elim{\lor} आणि
\Elim{\land} अनुमानांचा
क्रम बदलू (“क्रमपरिवर्तन”
करू) शकतो. म्हणजे~$\delta$
मधील उपसिद्धता~$\delta_1$
पुढील रचनेने बदलतो:
\[
\AxiomC{}
\RightLabel{$\delta_2$}
\DeduceC{$!D \lor !E$}
\AxiomC{$\Discharge{!D}{x}$}
\RightLabel{$\delta_3$}
\DeduceC{$!B \land !C$}
\RightLabel{\Elim{\land}[1]}
\UnaryInfC{$!B$}
\AxiomC{$\Discharge{!E}{y}$}
\RightLabel{$\delta_4$}
\DeduceC{$!B \land !C$}
\RightLabel{\Elim{\land}[1]}
\UnaryInfC{$!B$}
\DischargeRule{\Elim{\lor}}{x\,y}
\TrinaryInfC{$!B$}
\DisplayProof
\]
आधी~\Elim{\lor} च्या
खाली~\Elim{\land}
च्या आधारविधानात संपणारे
कट आता~\Elim{\lor}
च्या गौण आधारविधानांच्या
वर असलेल्या~\Elim{\land}
अनुमानांपैकी एका
आधारविधानात संपतात.
म्हणजे प्रत्येकाची लांबी
~$1$ ने कमी झाली.

याखेरीज~$\delta_2$,
$\delta_3$ किंवा~$\delta_4$
यांपैकी एखाद्यात पूर्णपणे
असलेला कट, किंवा~$\delta$
मधील~$\delta_1$
बाहेर पूर्णपणे असलेला कट,
बदलत नाही. त्यात तीच
!!{formula}s आणि तीच
अनुमाने असतात; म्हणून
विशेषतः त्याची कोटी व
लांबी~$\delta$ मधील
अनुरूप कटाइतकीच असते.
त्यामुळे नव्या !!{proof}
ची कट-कोटी वाढत नाही,
आणि नव्या !!{proof}
ची कट-लांबी
जुन्या !!{proof} च्या
कट-लांबीपेक्षा कमी होते.

\begin{prob}
क्रमपरिवर्तन रूपांतरणानंतर
किमान एका कटाची लांबी~$1$
ने कमी होते.
\Elim{\land} ला
एखाद्या~\Elim{\lor}
च्या पलीकडे हलवले,
तर प्रत्यक्षात दोन
कट-खंडांची लांबी कमी
होते; त्यामुळे या
प्रसंगी !!{proof} ची
कट-लांबी किमान~$2$
ने कमी होते.
अशा एकाच क्रमपरिवर्तनाने
!!a{proof} ची कट-लांबी
$2$ पेक्षा अधिक कमी
होऊ शकते का?
कटाची \emph{कोटी}
कमी होऊ शकते का?
\end{prob}

!!a{elimination} नियम
\Elim{\lor} किंवा
\Elim{\lexists} नियमाच्या
वर हलवण्याच्या प्रक्रियेला
\emph{क्रमपरिवर्तन रूपांतरण}
म्हणतात. काही नियम-जोड्यांचे
नमुने~\olref{tab:permutation-conversions}
मध्ये दिले आहेत.
(उरलेले सरावासाठी सोडले आहेत.)

% \begin{table}
% \begin{multline*}
% \shoveleft{\AxiomC{$!D \lor !E$}
% \AxiomC{$\Discharge{!D}{x}$}
% \shortDeduce
% \DeduceC{$!B \land !C$}
% \AxiomC{$\Discharge{!E}{y}$}
% \shortDeduce
% \DeduceC{$!B \land !C$}
% \DischargeRule{\Elim{\lor}}{x\,y}
% \TrinaryInfC{$!B \land !C$}
% \RightLabel{\Elim{\land}[1]}
% \UnaryInfC{$!B$}
% \DisplayProof}\\
% \shoveright{\longrightarrow\
% \AxiomC{$!D \lor !E$}
% \AxiomC{$\Discharge{!D}{x}$}
% \shortDeduce
% \DeduceC{$!B \land !C$}
% \RightLabel{\Elim{\land}[1]}
% \UnaryInfC{$!B$}
% \AxiomC{$\Discharge{!E}{y}$}
% \shortDeduce
% \DeduceC{$!B \land !C$}
% \RightLabel{\Elim{\land}[1]}
% \UnaryInfC{$!B$}
% \DischargeRule{\Elim{\lor}}{x\,y}
% \TrinaryInfC{$!B$}
% \DisplayProof}
% \\
% \shoveleft{\AxiomC{$!D \lor !E$}
% \AxiomC{$\Discharge{!D}{x}$}
% \shortDeduce
% \DeduceC{$!B \lif !C$}
% \AxiomC{$\Discharge{!E}{y}$}
% \shortDeduce
% \DeduceC{$!B \lif !C$}
% \DischargeRule{\Elim{\lor}}{x\,y}
% \TrinaryInfC{$!B \lif !C$}
% \AxiomC{$!B$}
% \RightLabel{\Elim{\lif}}
% \BinaryInfC{$!C$}
% \DisplayProof}\\
% \shoveright{\longrightarrow\
% \AxiomC{$!D \lor !E$}
% \AxiomC{$\Discharge{!D}{x}$}
% \shortDeduce
% \DeduceC{$!B \lif !C$}
% \AxiomC{$!B$}
% \RightLabel{\Elim{\lif}}
% \BinaryInfC{$!C$}
% \AxiomC{$\Discharge{!E}{y}$}
% \shortDeduce
% \DeduceC{$!B \lif !C$}
% \AxiomC{$!B$}
% \RightLabel{\Elim{\lif}}
% \BinaryInfC{$!C$}
% \DischargeRule{\Elim{\lor}}{x\,y}
% \TrinaryInfC{$!C$}
% \DisplayProof}
% \\
% \shoveleft{\AxiomC{$!D \lor !E$}
% \AxiomC{$\Discharge{!D}{x}$}
% \shortDeduce
% \DeduceC{$!B \lif !C$}
% \AxiomC{$\Discharge{!E}{y}$}
% \shortDeduce
% \DeduceC{$!B \lif !C$}
% \DischargeRule{\Elim{\lor}}{x\,y}
% \TrinaryInfC{$!B \lif !C$}
% \AxiomC{$!B$}
% \RightLabel{\Elim{\lif}}
% \BinaryInfC{$!C$}
% \DisplayProof}\\
% \shoveright{\longrightarrow\
% \AxiomC{$!D \lor !E$}
% \AxiomC{$\Discharge{!D}{x}$}
% \shortDeduce
% \DeduceC{$!B \lif !C$}
% \AxiomC{$!B$}
% \RightLabel{\Elim{\lif}}
% \BinaryInfC{$!C$}
% \AxiomC{$\Discharge{!E}{y}$}
% \shortDeduce
% \DeduceC{$!B \lif !C$}
% \AxiomC{$!B$}
% \RightLabel{\Elim{\lif}}
% \BinaryInfC{$!C$}
% \DischargeRule{\Elim{\lor}}{x\,y}
% \TrinaryInfC{$!C$}
% \DisplayProof}
% \\
% \shoveleft{\AxiomC{$!D \lor !E$}
% \AxiomC{$\Discharge{!D}{x}$}
% \shortDeduce
% \DeduceC{$\lforall[x][!B(x)]$}
% \AxiomC{$\Discharge{!E}{y}$}
% \shortDeduce
% \DeduceC{$\lforall[x][!B(x)]$}
% \DischargeRule{\Elim{\lor}}{x\,y}
% \TrinaryInfC{$\lforall[x][!B(x)]$}
% \RightLabel{\Elim{\lforall}}
% \UnaryInfC{$!B(t)$}
% \DisplayProof}\\
% \shoveright{\longrightarrow\
% \AxiomC{$!D \lor !E$}
% \AxiomC{$\Discharge{!D}{x}$}
% \shortDeduce
% \DeduceC{$\lforall[x][!B(x)]$}
% \RightLabel{\Elim{\lforall}}
% \UnaryInfC{$!B(t)$}
% \AxiomC{$\Discharge{!E}{y}$}
% \shortDeduce
% \DeduceC{$\lforall[x][!B(x)]$}
% \RightLabel{\Elim{\lforall}}
% \UnaryInfC{$!B(t)$}
% \DischargeRule{\Elim{\lor}}{x\,y}
% \TrinaryInfC{$!B(t)$}
% \DisplayProof}
% \end{multline*}
% \caption{Some permutation conversions for \Log{N1i}.}
% \ollabel{tab:permutation-conversions}
% \end{table}

{\small
\begin{longtable}{@{}l@{}}
\AxiomC{$!D \lor !E$}
\AxiomC{$\Discharge{!D}{x}$}
\shortDeduce
\DeduceC{$!B \land !C$}
\AxiomC{$\Discharge{!E}{y}$}
\shortDeduce
\DeduceC{$!B \land !C$}
\DischargeRule{\Elim{\lor}}{x\,y}
\TrinaryInfC{$!B \land !C$}
\RightLabel{\Elim{\land}[1]}
\UnaryInfC{$!B$}
\DisplayProof
$\longrightarrow$\\[6ex]
\quad\AxiomC{$!D \lor !E$}
\AxiomC{$\Discharge{!D}{x}$}
\shortDeduce
\DeduceC{$!B \land !C$}
\RightLabel{\Elim{\land}[1]}
\UnaryInfC{$!B$}
\AxiomC{$\Discharge{!E}{y}$}
\shortDeduce
\DeduceC{$!B \land !C$}
\RightLabel{\Elim{\land}[1]}
\UnaryInfC{$!B$}
\DischargeRule{\Elim{\lor}}{x\,y}
\TrinaryInfC{$!B$}
\DisplayProof
\\[10ex]
\AxiomC{$!D \lor !E$}
\AxiomC{$\Discharge{!D}{x}$}
\shortDeduce
\DeduceC{$!B \lif !C$}
\AxiomC{$\Discharge{!E}{y}$}
\shortDeduce
\DeduceC{$!B \lif !C$}
\DischargeRule{\Elim{\lor}}{x\,y}
\TrinaryInfC{$!B \lif !C$}
\AxiomC{$!B$}
\RightLabel{\Elim{\lif}}
\BinaryInfC{$!C$}
\DisplayProof
$\longrightarrow$\\[6ex]
\AxiomC{$!D \lor !E$}
\AxiomC{$\Discharge{!D}{x}$}
\shortDeduce
\DeduceC{$!B \lif !C$}
\AxiomC{$!B$}
\RightLabel{\Elim{\lif}}
\BinaryInfC{$!C$}
\AxiomC{$\Discharge{!E}{y}$}
\shortDeduce
\DeduceC{$!B \lif !C$}
\AxiomC{$!B$}
\RightLabel{\Elim{\lif}}
\BinaryInfC{$!C$}
\DischargeRule{\Elim{\lor}}{x\,y}
\TrinaryInfC{$!C$}
\DisplayProof
\\[10ex]
\AxiomC{$!D \lor !E$}
\AxiomC{$\Discharge{!D}{x}$}
\shortDeduce
\DeduceC{$!B \lif !C$}
\AxiomC{$\Discharge{!E}{y}$}
\shortDeduce
\DeduceC{$!B \lif !C$}
\DischargeRule{\Elim{\lor}}{x\,y}
\TrinaryInfC{$!B \lif !C$}
\AxiomC{$!B$}
\RightLabel{\Elim{\lif}}
\BinaryInfC{$!C$}
\DisplayProof
$\longrightarrow$\\[6ex]
\AxiomC{$!D \lor !E$}
\AxiomC{$\Discharge{!D}{x}$}
\shortDeduce
\DeduceC{$!B \lif !C$}
\AxiomC{$!B$}
\RightLabel{\Elim{\lif}}
\BinaryInfC{$!C$}
\AxiomC{$\Discharge{!E}{y}$}
\shortDeduce
\DeduceC{$!B \lif !C$}
\AxiomC{$!B$}
\RightLabel{\Elim{\lif}}
\BinaryInfC{$!C$}
\DischargeRule{\Elim{\lor}}{x\,y}
\TrinaryInfC{$!C$}
\DisplayProof
\\[10ex]
\AxiomC{$!D \lor !E$}
\AxiomC{$\Discharge{!D}{x}$}
\shortDeduce
\DeduceC{$\lforall[x][!B(x)]$}
\AxiomC{$\Discharge{!E}{y}$}
\shortDeduce
\DeduceC{$\lforall[x][!B(x)]$}
\DischargeRule{\Elim{\lor}}{x\,y}
\TrinaryInfC{$\lforall[x][!B(x)]$}
\RightLabel{\Elim{\lforall}}
\UnaryInfC{$!B(t)$}
\DisplayProof
\quad$\longrightarrow$\\[8ex]
\quad\AxiomC{$!D \lor !E$}
\AxiomC{$\Discharge{!D}{x}$}
\shortDeduce
\DeduceC{$\lforall[x][!B(x)]$}
\RightLabel{\Elim{\lforall}}
\UnaryInfC{$!B(t)$}
\AxiomC{$\Discharge{!E}{y}$}
\shortDeduce
\DeduceC{$\lforall[x][!B(x)]$}
\RightLabel{\Elim{\lforall}}
\UnaryInfC{$!B(t)$}
\DischargeRule{\Elim{\lor}}{x\,y}
\TrinaryInfC{$!B(t)$}
\DisplayProof
\\
\caption{\Log{N1i} साठी काही क्रमपरिवर्तन रूपांतरणे.}
\ollabel{tab:permutation-conversions}
\end{longtable}
}

\begin{prob}
\Elim{\lor} नंतर
\Elim{\lor} किंवा
\Elim{\lnot} येत असेल,
तर क्रमपरिवर्तन
रूपांतरणे द्या.
\end{prob}

\begin{prob}
\Elim{\lexists} नंतर
\Elim{\lexists} येत असेल,
तर क्रमपरिवर्तन
रूपांतरणे द्या.
दुसऱ्या प्रसंगी
आयगेन चराची कोणतीही
अट का मोडत नाही,
हे स्पष्ट करा.
\end{prob}

क्रमपरिवर्तन रूपांतरण
लावताना आयगेन चराची
अट मोडत नाही ना,
हे तपासले पाहिजे.
पुढील उदाहरण घ्या:
\[
\AxiomC{}
\RightLabel{$\delta_2$}
\DeduceC{$!D \lor !E$}
\AxiomC{$\Discharge{!D}{x}$}
\RightLabel{$\delta_3$}
\DeduceC{$\lexists[x][!B(x)]$}
\AxiomC{$\Discharge{!E}{y}$}
\RightLabel{$\delta_4$}
\DeduceC{$\lexists[x][!B(x)]$}
\DischargeRule{\Elim{\lor}}{x\,y}
\TrinaryInfC{$\lexists[x][!B(x)]$}
\AxiomC{$\Discharge{!B(c)}{z}$}
\RightLabel{$\delta_5$}
\DeduceC{$!C$}
\DischargeRule{\Elim{\lexists}}{z}
\BinaryInfC{$!C$}
\DisplayProof
\]
याच्या जागी पुढील
रचना बसवली, तर
\[
\AxiomC{}
\RightLabel{$\delta_2$}
\DeduceC{$!D \lor !E$}
\AxiomC{$\Discharge{!D}{x}$}
\RightLabel{$\delta_3$}
\DeduceC{$\lexists[x][!B(x)]$}
\AxiomC{$\Discharge{!B(c)}{z}$}
\RightLabel{$\delta_5$}
\DeduceC{$!C$}
\DischargeRule{\Elim{\lexists}}{z}
\BinaryInfC{$!C$}
\AxiomC{$\Discharge{!E}{y}$}
\RightLabel{$\delta_4$}
\DeduceC{$\lexists[x][!B(x)]$}
\AxiomC{$\Discharge{!B(c)}{z}$}
\RightLabel{$\delta_5$}
\DeduceC{$!C$}
\DischargeRule{\Elim{\lexists}}{z}
\BinaryInfC{$!C$}
\DischargeRule{\Elim{\lor}}{x\,y}
\TrinaryInfC{$!C$}
\DisplayProof
\]
आयगेन चर~$c$ हा,
उदाहरणार्थ,~$!D$
मध्ये येईल का,
अशी शंका येऊ शकते.
मूळ !!{proof} मध्ये
गृहीतक~$!D$ हे
मूळ~$\Elim{\lexists}$
च्या वर मुक्त होते.
म्हणून ते
\Elim{\lexists}
अनुमानाच्या प्रमुख
आधारविधानाच्या वर
उघड्या गृहीतकात
येत नाही; त्यामुळे
तेथे आयगेन चराची अट
पूर्ण होते. पण नव्या
!!{proof} मध्ये~$!D$
हे~\Elim{\lexists}
अनुमानांनंतरच मुक्त
होते, म्हणून ती अट
मोडेल असे वाटते.
तथापि, आपल्या
!!{proof}s नियमित
आहेत हे लक्षात घ्या:
प्रत्येक आयगेन चर
एकाच अनुमानाचा
आयगेन चर असतो आणि
तो !!{proof} मध्ये
फक्त~\Elim{\lexists}
च्या गौण आधारविधानाच्या
वरच येतो. विशेषतः
$c$ हा~$\delta_3$
किंवा~$\delta_4$
मध्ये येऊ शकत नाही.

या उदाहरणातून आणखी
एक गोष्ट दिसते.
नव्या !!{proof}
मध्ये~$\delta_5$
च्या दोन प्रती आहेत.
म्हणून~$\delta_5$
मध्ये महत्तम कट असेल,
तर आपण कट-लांबी
\emph{कमी केलेली नाही}!
त्यामुळे~$\delta_5$
मध्ये महत्तम कट नाही
हे माहीत असेल तेव्हाच
क्रमपरिवर्तन रूपांतरण
लावण्याची काळजी
घ्यावी लागते.
नेहमी \emph{सर्वात वरचा}
महत्तम कट निवडून
हे साधतो: म्हणजे असा
महत्तम कट की त्याच्या
उप-!!{proof} मध्ये
(येथे~$\delta_1$)
उप-!!{proof} च्या
शेवटच्या अनुमानाच्या
वर संपणारा दुसरा
महत्तम कट नसतो.
यामुळे~$\delta_5$
मध्ये, किंबहुना
$\delta_2$, $\delta_3$
किंवा~$\delta_4$
मध्येही, इतर महत्तम
कट असण्याची शक्यता
नाहीशी होते.

\end{document}
